% problem-000373 1 有机酸碱性
\section*{1. 有机酸碱性}
\textit{下列为分问。}

\subsection*{1-1}
2011年是国际化学年，是居⾥夫⼈获得诺⻉尔化学奖100周年。居⾥夫⼈发现的两种化学元素的元素符号和中⽂名称分别是 和

\subsection*{1-2}
向 $\mathrm { T i O S O _ { 4 } }$ ⽔溶液中加⼊锌粒，反应后溶液变为紫⾊。在清夜中滴加适量的 $\mathrm { C u C l _ { 2 } }$ ⽔溶液，产⽣⽩⾊沉淀。⽣成⽩⾊沉淀的离⼦⽅程式是 ；继续滴加 $\mathrm { C u C l _ { 2 } }$ ⽔溶液，⽩⾊沉淀消失，其离⼦⽅程式是

\subsection*{1-3}
20世纪60年代维也纳⼤学V.Gutmann研究⼩组报道，三原⼦分⼦A可由 $\mathrm { S F _ { 4 } }$ 和 $\mathrm { N H _ { 3 } }$ 反应合成；A被$\mathrm { A g F } _ { 2 }$ 氧化得到沸点为为 $2 7 ~ ^ { \circ } \mathrm { C }$ 的三元化合物 $\mathbf { B } _ { \circ }$ 。A和B分⼦中的中⼼原⼦与同种端位原⼦的核间距⼏乎相等；B分⼦有⼀根三种轴和3个镜⾯。画出A和B的结构式（明确⽰出单键和重键，不在纸⾯上的键⽤楔形键表${ \overline { { \overline { { \prime } } } \vec { \jmath } } } ,$ ⾮键合电⼦不必标出）。

\subsection*{1-4}
画出 $\mathrm { A l _ { 2 } ( n \mathrm { - C _ { 4 } H _ { 9 } ) _ { 4 } H _ { 2 } } }$ 和 $\mathrm { M g [ A l ( C H _ { 3 } ) _ { 4 } ] } ;$ _{2}的结构简式。

\subsection*{1-5}
已知 $1 E ^ { \Theta } ( \mathrm { F e O _ { 4 } } ^ { 2 - } / \mathrm { F e } ^ { 3 + } ) = 2 . 2 0 \mathrm { V } , E ^ { \Theta } ( \mathrm { F e O _ { 4 } } ^ { 2 - } / \mathrm { F e } ( \mathrm { O H } ) _ { 3 } ) = 0 . 7 2 \mathrm { V } _ { \circ }$

\subsection*{1-5}
-1 写出氯⽓和三氯化铁反应形成⾼铁酸根的离⼦⽅程 $\updownarrow \updownarrow _ { \updownarrow }$

\subsection*{1-5}
-2 写出⾼铁酸钾在酸性⽔溶液中分解的离⼦⽅程式。

\subsection*{1-5}
-3 ⽤⾼铁酸钾与镁等组成碱性电池，写出该电池的电极反应。
